% Options for packages loaded elsewhere
\PassOptionsToPackage{unicode}{hyperref}
\PassOptionsToPackage{hyphens}{url}
\documentclass[
  ignorenonframetext,
]{beamer}
\newif\ifbibliography
\usepackage{pgfpages}
\setbeamertemplate{caption}[numbered]
\setbeamertemplate{caption label separator}{: }
\setbeamercolor{caption name}{fg=normal text.fg}
\beamertemplatenavigationsymbolsempty
% remove section numbering
\setbeamertemplate{part page}{
  \centering
  \begin{beamercolorbox}[sep=16pt,center]{part title}
    \usebeamerfont{part title}\insertpart\par
  \end{beamercolorbox}
}
\setbeamertemplate{section page}{
  \centering
  \begin{beamercolorbox}[sep=12pt,center]{section title}
    \usebeamerfont{section title}\insertsection\par
  \end{beamercolorbox}
}
\setbeamertemplate{subsection page}{
  \centering
  \begin{beamercolorbox}[sep=8pt,center]{subsection title}
    \usebeamerfont{subsection title}\insertsubsection\par
  \end{beamercolorbox}
}
% Prevent slide breaks in the middle of a paragraph
\widowpenalties 1 10000
\raggedbottom
\AtBeginPart{
  \frame{\partpage}
}
\AtBeginSection{
  \ifbibliography
  \else
    \frame{\sectionpage}
  \fi
}
\AtBeginSubsection{
  \frame{\subsectionpage}
}
\usepackage{iftex}
\ifPDFTeX
  \usepackage[T1]{fontenc}
  \usepackage[utf8]{inputenc}
  \usepackage{textcomp} % provide euro and other symbols
\else % if luatex or xetex
  \usepackage{unicode-math} % this also loads fontspec
  \defaultfontfeatures{Scale=MatchLowercase}
  \defaultfontfeatures[\rmfamily]{Ligatures=TeX,Scale=1}
\fi
\usepackage{lmodern}
\ifPDFTeX\else
  % xetex/luatex font selection
\fi
% Use upquote if available, for straight quotes in verbatim environments
\IfFileExists{upquote.sty}{\usepackage{upquote}}{}
\IfFileExists{microtype.sty}{% use microtype if available
  \usepackage[]{microtype}
  \UseMicrotypeSet[protrusion]{basicmath} % disable protrusion for tt fonts
}{}
\makeatletter
\@ifundefined{KOMAClassName}{% if non-KOMA class
  \IfFileExists{parskip.sty}{%
    \usepackage{parskip}
  }{% else
    \setlength{\parindent}{0pt}
    \setlength{\parskip}{6pt plus 2pt minus 1pt}}
}{% if KOMA class
  \KOMAoptions{parskip=half}}
\makeatother
\usepackage{color}
\usepackage{fancyvrb}
\newcommand{\VerbBar}{|}
\newcommand{\VERB}{\Verb[commandchars=\\\{\}]}
\DefineVerbatimEnvironment{Highlighting}{Verbatim}{commandchars=\\\{\}}
% Add ',fontsize=\small' for more characters per line
\newenvironment{Shaded}{}{}
\newcommand{\AlertTok}[1]{\textcolor[rgb]{1.00,0.00,0.00}{\textbf{#1}}}
\newcommand{\AnnotationTok}[1]{\textcolor[rgb]{0.38,0.63,0.69}{\textbf{\textit{#1}}}}
\newcommand{\AttributeTok}[1]{\textcolor[rgb]{0.49,0.56,0.16}{#1}}
\newcommand{\BaseNTok}[1]{\textcolor[rgb]{0.25,0.63,0.44}{#1}}
\newcommand{\BuiltInTok}[1]{\textcolor[rgb]{0.00,0.50,0.00}{#1}}
\newcommand{\CharTok}[1]{\textcolor[rgb]{0.25,0.44,0.63}{#1}}
\newcommand{\CommentTok}[1]{\textcolor[rgb]{0.38,0.63,0.69}{\textit{#1}}}
\newcommand{\CommentVarTok}[1]{\textcolor[rgb]{0.38,0.63,0.69}{\textbf{\textit{#1}}}}
\newcommand{\ConstantTok}[1]{\textcolor[rgb]{0.53,0.00,0.00}{#1}}
\newcommand{\ControlFlowTok}[1]{\textcolor[rgb]{0.00,0.44,0.13}{\textbf{#1}}}
\newcommand{\DataTypeTok}[1]{\textcolor[rgb]{0.56,0.13,0.00}{#1}}
\newcommand{\DecValTok}[1]{\textcolor[rgb]{0.25,0.63,0.44}{#1}}
\newcommand{\DocumentationTok}[1]{\textcolor[rgb]{0.73,0.13,0.13}{\textit{#1}}}
\newcommand{\ErrorTok}[1]{\textcolor[rgb]{1.00,0.00,0.00}{\textbf{#1}}}
\newcommand{\ExtensionTok}[1]{#1}
\newcommand{\FloatTok}[1]{\textcolor[rgb]{0.25,0.63,0.44}{#1}}
\newcommand{\FunctionTok}[1]{\textcolor[rgb]{0.02,0.16,0.49}{#1}}
\newcommand{\ImportTok}[1]{\textcolor[rgb]{0.00,0.50,0.00}{\textbf{#1}}}
\newcommand{\InformationTok}[1]{\textcolor[rgb]{0.38,0.63,0.69}{\textbf{\textit{#1}}}}
\newcommand{\KeywordTok}[1]{\textcolor[rgb]{0.00,0.44,0.13}{\textbf{#1}}}
\newcommand{\NormalTok}[1]{#1}
\newcommand{\OperatorTok}[1]{\textcolor[rgb]{0.40,0.40,0.40}{#1}}
\newcommand{\OtherTok}[1]{\textcolor[rgb]{0.00,0.44,0.13}{#1}}
\newcommand{\PreprocessorTok}[1]{\textcolor[rgb]{0.74,0.48,0.00}{#1}}
\newcommand{\RegionMarkerTok}[1]{#1}
\newcommand{\SpecialCharTok}[1]{\textcolor[rgb]{0.25,0.44,0.63}{#1}}
\newcommand{\SpecialStringTok}[1]{\textcolor[rgb]{0.73,0.40,0.53}{#1}}
\newcommand{\StringTok}[1]{\textcolor[rgb]{0.25,0.44,0.63}{#1}}
\newcommand{\VariableTok}[1]{\textcolor[rgb]{0.10,0.09,0.49}{#1}}
\newcommand{\VerbatimStringTok}[1]{\textcolor[rgb]{0.25,0.44,0.63}{#1}}
\newcommand{\WarningTok}[1]{\textcolor[rgb]{0.38,0.63,0.69}{\textbf{\textit{#1}}}}
\setlength{\emergencystretch}{3em} % prevent overfull lines
\providecommand{\tightlist}{%
  \setlength{\itemsep}{0pt}\setlength{\parskip}{0pt}}
% Slide layout defaults for warning-clean scientific decks.
\usepackage{etoolbox}
\IfFileExists{xurl.sty}{\usepackage{xurl}}{}
\IfFileExists{seqsplit.sty}{\usepackage{seqsplit}}{\newcommand{\seqsplit}[1]{#1}}
\protected\def\breakseq#1{\seqsplit{#1}}
\protected\def\breaktt#1{\begingroup\ttfamily\seqsplit{#1}\endgroup}
\setlength{\emergencystretch}{6em}
\tolerance=5000
\hbadness=10000
\hfuzz=1pt
\setlength{\tabcolsep}{2pt}
\AtBeginEnvironment{longtable}{\tiny\renewcommand{\arraystretch}{0.86}\setlength{\tabcolsep}{1pt}}
\AtBeginEnvironment{tabular}{\tiny\renewcommand{\arraystretch}{0.86}\setlength{\tabcolsep}{1pt}}

% Natbib and cross-reference fallbacks — slides don't load natbib
% or cleveref, but combined-PDF manuscript prose may emit these
% commands. The fallback renders citations as a bracketed key list
% and cross-references as detokenized labels so slides don't fail on
% undefined control sequences or raw underscores. \providecommand is
% a no-op if packages load later.
\providecommand{\citep}[1]{[#1]}
\providecommand{\citet}[1]{#1}
\providecommand{\citealp}[1]{#1}
\providecommand{\citeauthor}[1]{#1}
\providecommand{\citeyear}[1]{#1}
\providecommand{\cref}[1]{\texttt{\detokenize{#1}}}
\providecommand{\Cref}[1]{\texttt{\detokenize{#1}}}
\usepackage{bookmark}
\IfFileExists{xurl.sty}{\usepackage{xurl}}{} % add URL line breaks if available
\urlstyle{same}
% fallback for those not using the hyperref driver hyperxmp:
\makeatletter
\@ifundefined{xmpquote}{\newcommand{\xmpquote}[1]{#1}}{}
\makeatother
\hypersetup{
  hidelinks,
  pdfcreator={LaTeX via pandoc}}

\author{\texorpdfstring{}{}}
\date{}

\begin{document}

\section{Deep Search}\label{sec:deep_search}

\begin{frame}[allowframebreaks]{Deep Search}
The deep-search workflow extends the standard literature pipeline along
three axes: \textbf{breadth} (multi-keyword fan-out), \textbf{depth}
(full enrichment of every paper, abstract + PDF fulltext), and
\textbf{archival quality} (per-paper multi-section LLM reading notes
saved as standalone markdown). It is invoked from the same configuration
file but a separate orchestrator script.
\end{frame}

\begin{frame}[fragile,allowframebreaks]{Configuration}
\protect\phantomsection\label{configuration}
The \texttt{deep\_search:} block in \breaktt{manuscript/config.yaml}
controls the run. The most-used knobs:

\begin{itemize}
\tightlist
\item
  \texttt{keywords} --- list of free-text queries (each becomes one
  \texttt{SearchQuery}).
\item
  \breaktt{max\_results\_per\_keyword} --- per-keyword cap (100 by
  default; honoured by the aggregator after dedup).
\item
  \texttt{sources} --- backend list, same vocabulary as the standard
  pipeline (\texttt{arxiv}, \texttt{crossref}, \texttt{local},
  \texttt{paperclip}).
\item
  \texttt{fetch\_abstracts} / \texttt{fetch\_fulltext} --- when both are
  \texttt{true}, every returned paper has its abstract fetched (arXiv
  export API) and (where a PDF URL is available) its fulltext extracted
  via \texttt{pypdf}. Cached on disk under \breaktt{output/cache/abs/}
  and \breaktt{output/cache/pdf/}.
\item
  \texttt{llm\_per\_paper} --- when \texttt{true} and Ollama is
  reachable, each paper gets a multi-section markdown reading note
  generated by the local LLM. When \texttt{false}, or when the LLM stack
  is genuinely unreachable at runtime, the per-paper note is written
  with only the abstract / fulltext-excerpt sections; the synthesis
  section is omitted entirely (no placeholder text). The composer's
  Supplemental S1 also drops the per-paper synthesis subsection in that
  case rather than emitting empty rows.
\item
  \breaktt{write\_unified\_bibtex} / \breaktt{unified\_bibtex\_path} ---
  when \texttt{true}, a deduplicated, collision-suffixed BibTeX file is
  written under \breaktt{manuscript/references\_deep.bib} so it can be
  cited from the manuscript via Pandoc \texttt{\breakseq{[@key]}} syntax
  exactly the same way as the hand-curated \texttt{references.bib}.
\end{itemize}
\end{frame}

\begin{frame}[fragile,allowframebreaks]{Pipeline shape}
\protect\phantomsection\label{pipeline-shape}
The deep-search pipeline reads the \texttt{deep\_search:} block from
\breaktt{manuscript/config.yaml}, runs one \texttt{SearchQuery} per
keyword (capped at \breaktt{max\_results\_per\_keyword}), aggregates the
per-backend results through \breaktt{LiteratureClient}, enriches every
paper via \texttt{AbstractFetcher} and \texttt{FulltextFetcher},
generates collision-free citation keys with
\breaktt{paper\_to\_bibentry}, and (when \breaktt{llm\_per\_paper:\ true}
and Ollama is reachable) calls the local LLM with \texttt{DEEP\_PROMPT}
to produce a seven-section reading note per paper. The per-keyword
outputs are then merged by \texttt{merge\_papers} to form a deduplicated
aggregate roster which is written to
\breaktt{manuscript/references\_deep.bib},
\breaktt{output/deep\_search/aggregate.json}, and
\breaktt{output/deep\_search/aggregate\_report.md}. A Mermaid flowchart
in this subsection renders the same flow visually in the HTML build.

\begin{Shaded}
\begin{Highlighting}[]
\NormalTok{flowchart TB}
\NormalTok{    CFG[deep\_search block of config.yaml]}
\NormalTok{    CFG {-}{-}\textgreater{} KW\{\{For each keyword\}\}}
\NormalTok{    KW {-}{-}\textgreater{} SEARCH[SearchQuery · max\_results\_per\_keyword]}
\NormalTok{    SEARCH {-}{-}\textgreater{} CLIENT[LiteratureClient\textless{}br/\textgreater{}aggregates arxiv · crossref · local · paperclip]}
\NormalTok{    CLIENT {-}{-}\textgreater{} ENRICH[AbstractFetcher + FulltextFetcher\textless{}br/\textgreater{}cached on disk]}
\NormalTok{    ENRICH {-}{-}\textgreater{} KEYS[paper\_to\_bibentry\textless{}br/\textgreater{}collision{-}free keys]}
\NormalTok{    KEYS {-}{-}\textgreater{} LLMQ\{LLM per{-}paper?\}}
\NormalTok{    LLMQ {-}{-} yes {-}{-}\textgreater{} LLM[Ollama LLMClient\textless{}br/\textgreater{}DEEP\_PROMPT · 7 sections]}
\NormalTok{    LLMQ {-}{-} no {-}{-}\textgreater{} SKIP[skip synthesis · no placeholder]}
\NormalTok{    LLM {-}{-}\textgreater{} NOTE[output/deep\_search/\&lt;slug\&gt;/per\_paper/\&lt;safe\_id\&gt;.md]}
\NormalTok{    SKIP {-}{-}\textgreater{} NOTE}
\NormalTok{    NOTE {-}{-}\textgreater{} KR[KeywordResult]}

\NormalTok{    KR {-}{-}\textgreater{} AGG[merge\_papers across all keywords]}
\NormalTok{    AGG {-}{-}\textgreater{} BIB[unified BibTeX\textless{}br/\textgreater{}references\_deep.bib]}
\NormalTok{    AGG {-}{-}\textgreater{} AJSON[aggregate.json]}
\NormalTok{    AGG {-}{-}\textgreater{} AMD[aggregate\_report.md]}

\NormalTok{    classDef cfg fill:\#0f172a,stroke:\#0f172a,color:\#fff}
\NormalTok{    classDef proc fill:\#1e3a8a,stroke:\#0f172a,color:\#fff}
\NormalTok{    classDef llm fill:\#7c2d12,stroke:\#0f172a,color:\#fff}
\NormalTok{    classDef out fill:\#0f766e,stroke:\#0f172a,color:\#fff}
\NormalTok{    class CFG,KW,LLMQ cfg}
\NormalTok{    class SEARCH,CLIENT,ENRICH,KEYS,KR,AGG proc}
\NormalTok{    class LLM,SKIP llm}
\NormalTok{    class NOTE,BIB,AJSON,AMD out}
\end{Highlighting}
\end{Shaded}
\end{frame}

\begin{frame}[fragile,allowframebreaks]{LLM prompt}
\protect\phantomsection\label{llm-prompt}
The deep-search prompt is richer than the standard
\breaktt{synthesis.PROMPT\_PER\_PAPER}. It produces a self-contained
reading note suitable for archival without re-reading the paper:

\begin{verbatim}
## Contribution
One paragraph stating the paper's central claim and why it is novel.

## Method
2-4 bullets describing the technical approach.

## Evidence
2-4 bullets describing experiments / proofs.

## Limitations
1-3 bullets covering caveats and what the paper does NOT address.

## Connections
1-3 bullets relating this paper to other work in the field, citing only
papers explicitly named in the input.

## Significance for {keyword}
One paragraph explaining why this paper matters for the keyword.

## Tags
5-10 lowercase keywords.
\end{verbatim}

The LLM is duck-typed as \texttt{Callable{[}{[}str{]},\ str{]}} so tests
pass a deterministic local callable that returns real, well-formed
reading-note text; runtime callers wrap an
\breaktt{infrastructure.llm.LLMClient} with
\breaktt{seed=42,\ temperature=0.0} for reproducibility. When the LLM
stack is unreachable, callers pass \texttt{None} and the per-paper
synthesis stage is skipped entirely --- no placeholder text is ever
written into the archive.
\end{frame}

\begin{frame}[fragile,allowframebreaks]{On-disk layout}
\protect\phantomsection\label{on-disk-layout}
A successful deep-search run writes the following artefacts (paths shown
relative to the project root); a Mermaid tree in this subsection renders
the same hierarchy in the HTML build:

\begin{itemize}
\tightlist
\item
  \breaktt{output/deep\_search/aggregate.json} --- every unique paper
  across keywords.
\item
  \breaktt{output/deep\_search/aggregate\_report.md} --- cross-keyword
  markdown summary.
\item
  \breaktt{output/deep\_search/run\_summary.json} --- one-line metadata
  for the run.
\item
  \texttt{output/deep\_search/\textless{}keyword\_slug\textgreater{}/papers.json}
  --- enriched per-keyword paper list.
\item
  \texttt{output/deep\_search/\textless{}keyword\_slug\textgreater{}/reading\_report.md}
  --- per-keyword markdown summary.
\item
  \texttt{output/deep\_search/\textless{}keyword\_slug\textgreater{}/per\_paper/\textless{}safe\_id\textgreater{}.md}
  --- one reading note per paper.
\item
  \breaktt{manuscript/references\_deep.bib} --- auto-generated,
  deduplicated unified BibTeX file.
\end{itemize}

\begin{Shaded}
\begin{Highlighting}[]
\NormalTok{flowchart TB}
\NormalTok{    P["projects/templates/template\_search\_project/"]}
\NormalTok{    P {-}{-}\textgreater{} OUT["output/deep\_search/"]}
\NormalTok{    P {-}{-}\textgreater{} BIB[manuscript/references\_deep.bib\textless{}br/\textgreater{}auto{-}generated · deduplicated]}

\NormalTok{    OUT {-}{-}\textgreater{} AGG[aggregate.json\textless{}br/\textgreater{}every unique paper]}
\NormalTok{    OUT {-}{-}\textgreater{} AMD[aggregate\_report.md\textless{}br/\textgreater{}cross{-}keyword summary]}
\NormalTok{    OUT {-}{-}\textgreater{} RS[run\_summary.json\textless{}br/\textgreater{}one{-}line metadata]}
\NormalTok{    OUT {-}{-}\textgreater{} KW1["\&lt;keyword\_slug\_1\&gt;/"]}
\NormalTok{    OUT {-}{-}\textgreater{} KW2["\&lt;keyword\_slug\_2\&gt;/"]}
\NormalTok{    OUT {-}{-}\textgreater{} KW3["\&lt;keyword\_slug\_3\&gt;/"]}

\NormalTok{    KW1 {-}{-}\textgreater{} KW1\_P[papers.json]}
\NormalTok{    KW1 {-}{-}\textgreater{} KW1\_R[reading\_report.md]}
\NormalTok{    KW1 {-}{-}\textgreater{} KW1\_PP["per\_paper/\textless{}br/\textgreater{}SAFE\_ID.md per paper"]}

\NormalTok{    classDef d fill:\#0f172a,stroke:\#0f172a,color:\#fff}
\NormalTok{    classDef f fill:\#0f766e,stroke:\#0f172a,color:\#fff}
\NormalTok{    classDef bib fill:\#1e3a8a,stroke:\#0f172a,color:\#fff}
\NormalTok{    class P,OUT,KW1,KW2,KW3,KW1\_PP d}
\NormalTok{    class AGG,AMD,RS,KW1\_P,KW1\_R f}
\NormalTok{    class BIB bib}
\end{Highlighting}
\end{Shaded}
\end{frame}

\begin{frame}[fragile,allowframebreaks]{Determinism}
\protect\phantomsection\label{determinism}
Three caches make this workflow replayable:

\begin{enumerate}
\tightlist
\item
  \texttt{SearchCache}
  (\texttt{output/search/cache/search\_\textless{}hash\textgreater{}.json})
  --- keyed on canonical query identity per keyword.
\item
  Abstract cache
  (\texttt{output/cache/abs/\textless{}safe\_id\textgreater{}.txt}).
\item
  Fulltext cache
  (\texttt{output/cache/pdf/\textless{}safe\_id\textgreater{}.\{pdf,txt\}}).
\item
  LLM seed (\breaktt{deep\_search.llm\_seed:\ 42}) + temperature
  (\texttt{0.0}).
\end{enumerate}

Commit any subset of these caches to version control to freeze a run.
\end{frame}

\begin{frame}[fragile,allowframebreaks]{Paperclip backend}
\protect\phantomsection\label{paperclip-backend}
The \href{https://paperclip.gxl.ai}{Paperclip backend} is opt-in. Enable
it by:

\begin{enumerate}
\tightlist
\item
  Creating \breaktt{projects/templates/template\_search\_project/.env}
  (gitignored) with your \breaktt{PAPERCLIP\_API\_KEY=gxl\_…} (template
  at \texttt{.env.example}).
\item
  Including \texttt{paperclip} in \breaktt{deep\_search.sources}.
\end{enumerate}

The orchestrator scripts auto-load \texttt{.env} via the lightweight
\texttt{src.dotenv} module before constructing the backend list. The
\breaktt{PaperclipBackend} adapter mirrors the wire protocol of the
official \texttt{gxl\_paperclip} Python SDK: \texttt{POST\ /papers} with
an \texttt{X-API-Key} header and a JSON-RPC \texttt{tools/call}
envelope. We support both the modern \breaktt{structuredContent.papers}
response shape and the older text-only \texttt{content{[}{]}.text} shape
(best-effort regex extraction).

Failure-isolation: any per-backend error (including the migration-state
HTTP 405 the production endpoint may currently return) is captured into
\texttt{SearchResult.errors{[}paperclip{]}} and the run continues. See
\texttt{output/deep\_search/\textless{}keyword\textgreater{}/papers.json}
and the aggregate \breaktt{run\_summary.json} for the recorded errors.
\end{frame}

\begin{frame}[fragile,allowframebreaks]{CLI}
\protect\phantomsection\label{cli}
\begin{Shaded}
\begin{Highlighting}[]
\CommentTok{\# Run with everything from config.yaml (assumes deep\_search.enabled: true)}
\ExtensionTok{uv}\NormalTok{ run python projects/templates/template\_search\_project/scripts/run\_deep\_search.py}

\CommentTok{\# Force{-}enable without touching config.yaml}
\ExtensionTok{uv}\NormalTok{ run python projects/templates/template\_search\_project/scripts/run\_deep\_search.py }\AttributeTok{{-}{-}enable}

\CommentTok{\# Override keyword list at the command line}
\ExtensionTok{uv}\NormalTok{ run python projects/templates/template\_search\_project/scripts/run\_deep\_search.py }\DataTypeTok{\textbackslash{}}
    \AttributeTok{{-}{-}enable} \AttributeTok{{-}{-}keyword} \StringTok{"convex optimization"} \AttributeTok{{-}{-}keyword} \StringTok{"stochastic gradient descent"}

\CommentTok{\# Skip LLM stage even when config enables it}
\ExtensionTok{uv}\NormalTok{ run python projects/templates/template\_search\_project/scripts/run\_deep\_search.py }\AttributeTok{{-}{-}enable} \AttributeTok{{-}{-}no{-}llm}

\CommentTok{\# Bypass cache (writes still happen)}
\ExtensionTok{uv}\NormalTok{ run python projects/templates/template\_search\_project/scripts/run\_deep\_search.py }\AttributeTok{{-}{-}enable} \AttributeTok{{-}{-}no{-}cache}

\CommentTok{\# Local{-}corpus mode (offline · CI{-}friendly)}
\ExtensionTok{uv}\NormalTok{ run python projects/templates/template\_search\_project/scripts/run\_deep\_search.py }\DataTypeTok{\textbackslash{}}
    \AttributeTok{{-}{-}enable} \AttributeTok{{-}{-}corpus}\NormalTok{ projects/templates/template\_search\_project/data/corpus.json}
\end{Highlighting}
\end{Shaded}
\end{frame}

\end{document}
